% !TeX root = ../papers/nc-p3.tex
\section{文字交換の固定直線上の楕円曲線と正則性}
\label{sec:fixed-line-elliptic}
通常倍率、すなわち
$z\lambda(z^2-1)(\lambda^2z-1)(\lambda^2-1)(\lambda^2z^2-1)\ne0$
を取る。前節の固有直線を $L_\sigma$ ($\sigma=\pm1$) とし、その有限座標と補完を
\[
 A=\begin{pmatrix}a&1\\\sigma z&\sigma a\end{pmatrix},
 \qquad B=\begin{pmatrix}e&f\\\sigma zf&\sigma e\end{pmatrix}
\]
と書く。以下の点と開集合は標数零の幾何点について述べる。

\begin{theorem}\label{thm:fixed-line-elliptic}
各 $L_\sigma$ 上の符号付き対根の滑らかな完備化は楕円曲線である。
二曲線の $j$ 不変量は等しく、
\[
 j=256\frac{(1-k+k^2)^3}{k^2(1-k)^2},\qquad
 k=\frac{(\lambda^2-1)(\lambda^2z^2-1)}{(\lambda^2z-1)^2}.
\]
各直線はKummer四次面と横断的に四点で交わり、八点全てで両行列が可逆な真の重根を得る。
$(z,\lambda)=(3,2)$を含む倍率平面の非空開集合では、
八点全てが反対階数16の標準Koszul左右AS正則窓となる。
各始点は四つの中心四次式の多項式環上で階数256の自由加群である。
この開集合の各倍率で、各楕円曲線上の可逆窓は有限個の候補を除いて同じ性質を持つ。
\end{theorem}
\begin{proof}
$P(x)=px^2+mx+n$, $Q(y)=qy^2+hy+o$とし、
$F_\sigma=\lambda^2(z^2-1)^2(a^2-z)^2(m^2-4pn)$と置く。
$a$の四次式の最高次係数は $c=\lambda^2(z^2-1)^2$、定数項は $z^2c$ で、
その判別式は
\[
 2^{12}\lambda^4 z^6(\lambda^2-1)^2(\lambda^2 z^2-1)^2
 (\lambda^2z-1)^4(z-1)^4(z-\sigma)^{12}.
\]
従って四根は相異なる有限非零点であり、制限四次の単根なので面の滑らかな横断点である。
二直線は共通の射影点を持たない。二重被覆 $s^2=F_\sigma$ はRiemann--Hurwitzにより種数1、
無限遠に有理点を持つ。四次式の係数を降冪順に $A',B',C',D',E'$ と書き、
$I=12A'E'-3B'D'+(C')^2$を用いると $j=256I^3/\mathrm{disc}(F_\sigma)$。
$k=u/v$を上記の分子・分母で表すと
$I=16z^2(z-\sigma)^4(u^2-uv+v^2)$ となり、表示を得る。

$F_\sigma$と$a^2-z$の終結式は $16z^4(z-\sigma)^4(\lambda^2z-1)^4$。
重根公式 $x=-m/(2p),\ y=-h/(2q)$を補完式に入れた
$p,q,x,y,e,f,e^2-zf^2$ の各分子・分母についても、既約因子と
$F_\sigma$との終結式の因子は通常倍率で逆にした因子に限られる。
\path{computations/verify_fixed_line_elliptic_curves.py} がこれらの多項式恒等式を検算する。
ゆえに重根で元の可逆表示が定まる。

倍率 $(3,2)$ で原始化した二式は順に
\[
 16a^4-156a^3+555a^2-468a+144,\qquad 4a^4+339a^2+36.
\]
素数307ではそれぞれ $102,153,298,301$ と $80,119,188,227$ で単根に分裂する。
\path{computations/verify_fixed_branch_regular.py} は各剰余点の二行列から
四つの関係空間を作る。隣接階数はともに6、反対階数は16、
全始点で長さ4,5の次元は35,56、四次中心次元は文字順に $(3,2,0,0)$ である。
中心の最初の二文字次数から各二式を選ぶ。商の次元は全八点・四始点のいずれでも、
次数0から13の順に
\[
 1,4,10,20,31,40,44,40,31,20,10,4,1,0
\]
となる。

単根なので各剰余点は307進整数環に一意に持ち上がり、四根を両式とも尽くす。
各持上げで経路次元の上界は剰余体での値、下界は五次元対根族の正則一般点での値である。
従って長さ4,5の経路空間はこの離散付値環上で自由であり、中心方程式を同じ基底で比較できる。
一般点の中心からの下界5と剰余体の上界5が等しいので方程式は定階数で、
選んだ四式は中心に持ち上がる。商次元の上界と元の全次数の下界を合わせ、
定理\ref{thm:klein-rational-period-four}の自由基底と左右Koszul列の証明を適用する。
体拡大での完全性を降下でき、中心四式を選べる条件も核空間内の非空開条件である。
従って標数零の幾何点で主張が成り立つ。

長さ4,5の経路空間、中心方程式、商の長さ13までの階数条件は同じ議論で
各標本の開近傍に延びる。二つの四次式の根の軌跡は通常倍率上で有限エタールである。
八近傍の外の閉集合をこの有限射で倍率平面に写し、その閉像を除けば
$(3,2)$を残して八点を同時に扱える。固定倍率の滑らかな各被覆曲線にも
この正則な近傍が開集合として交わり、その補集合は有限である。
\end{proof}

\begin{proposition}\label{prop:fixed-line-defects}
各通常倍率で次の十二表示は全て真の相異なる対根の族に属し、長さ4で標準性を失う。
一行に同じ $a$ での二つの補完を記す。
\[
\begin{array}{cc|cc}
 \sigma&a&(e,f)&(e,f)\\ \hline
 +&\lambda z&(1,\lambda)&(z,\lambda z)\\
 +&1/\lambda&(z,1/\lambda)&(1,1/(\lambda z))\\
 -&\lambda z&(-1,\lambda)&(z,-\lambda z)\\
 -&-\lambda z&(1,\lambda)&(-z,-\lambda z)\\
 -&1/\lambda&(-z,1/\lambda)&(1,-1/(\lambda z))\\
 -&-1/\lambda&(z,1/\lambda)&(-1,-1/(\lambda z))
\end{array}
\]
\end{proposition}
\begin{proof}
$x=\sigma(a^2-zf^2)$, $y=\sigma(e^2-z)$で $P(x)=Q(y)=0$ かつ
$2px+m=-(2qy+h)\ne0$ となる。両行列式と $p,q,x,y$ も通常倍率では消えない。
前記の曲線の検算コードは十二表示について、関数体 $\mathbb Q(z,\lambda)$ 上で
四辺の六関係の全ての両側挿入を長さ4まで掃き出す。
正の四表示では全始点で次元36、負の八表示では始点を並べ替えて
$35,35,36,36$ である。通常倍率の特殊化でも関係は両切断で定階数6であり、
商次元は減少しない。標準窓の次元35に反する始点が常に残る。
\end{proof}
